\documentclass[11pt, oneside]{article} 
\usepackage{geometry}
\geometry{letterpaper} 
\usepackage{graphicx}
	
\usepackage{amssymb}
\usepackage{amsmath}
\usepackage{parskip}
\usepackage{color}
\usepackage{hyperref}

\graphicspath{{/Users/telliott/Dropbox/Github-math/figures/}}
% \begin{center} \includegraphics [scale=0.4] {gauss3.png} \end{center}

\title{Broken Chord}
\date{}

\begin{document}
\maketitle
\Large

%[my-super-duper-separator]

As mentioned previously, the theorem of the ``broken chord'' is ascribed to Archimedes, although his original work --- the \emph{Book of Circles} --- has been lost.  

It was analyzed in proofs collected by the Arabic mathematician Al Biruni in his \emph{Book on the Derivation of Chords in a Circle}.

Here is the general setup:
\begin{center} \includegraphics [scale=0.2] {BC_0.png} \end{center}

Let $A$ and $G$ be any two points on a circle, and let $D$ be equidistant from both, so that arc $AD$ is equal to arc $GD$.  Let $B$ be another point on the circle, lying between $G$ and $D$.

Drop the perpendicular from $D$ to $E$ on $BC$.  

The claim of the theorem is that $GB + BE = AE$.

I have come across even more proofs in recent years.  Recently, I came across a translation of Al Biruni's book, written in German by a historian named Heinrich Suter.

\url{http://www.jphogendijk.nl/biruni/Suter-Chords.pdf}

Here are some additional proofs, five from Suter and two from the web.

\subsection*{fourth proof:  Suter (i)}

This one is attributed to El-Sidjzi (972), but may actually have been known to Apollonius.

\url{https://www.researchgate.net/publication/341579803_FROM_THE_THEOREM_OF_THE_BROKEN_CHORD_TO_THE_BEGINNING_OF_TRIGONOMETRY}

This figure from Al Biruni (below) suggests a neat line of attack.

\begin{center} \includegraphics [scale=0.35] {Al_Biruni_5.png} \end{center}

Recall from our work on parallel chords in a circle, a result about the extensions from a rectangle in a circle.  In the figure below, $DE = CF$.

\begin{center} \includegraphics [scale=0.16] {rect_in_circle2.png} \end{center}

\emph{Proof}.  

Let $AB$ and $EF$ be two parallel chords in a circle with unequal lengths.  Draw the perpendiculars $AD$ and $BC$.  Then $ABCD$ is a rectangle.  

$\angle AFE = \angle BAF$ by alternate interior angles.  So arc $AE = $ arc $BF$.

It follows easily that $\triangle ADE \cong \triangle BCF$ by HL.

$\square$

\emph{Proof}.  Alternate.  Draw the radius which bisects $AB$.  Since $AB \parallel CD$, the same radius bisects $CD$ and it also bisects $EF$, since they are collinear.  If $M$ is the point of bisection on $CD$, then $ME = MF$ and the result follows easily by subtraction. $\square$

So now (changing notation):

\begin{center} \includegraphics [scale=0.18] {BC_i.png} \end{center}

\emph{Proof}.

As before, arc $GD = $ arc $AD$, and $DE \perp AEB$.

Erect the perpendicular to $DE$ at $D$ and find where it cuts the circle at $Z$.  This always forms a polygon since $BD < AD$.

Erect the perpendicular $ZH$. 

Since it has all right angles, $DZHE$ is a rectangle.

By our preliminary work, $\triangle DBE \cong \triangle ZAH$.

Thus arc $BD = $ arc $AZ$ and $BE = AH$.

Subtracting equals, arc $DZ = $ arc $GB$.

It follows that $GB = DZ = EH$.

Adding equals:

$GB + BE = EH + AH = AE$.

$\square$

\subsection*{fifth proof}

Here is another diagram from Al Biruni's book in Arabic.

\begin{center} \includegraphics [scale=0.35] {Al_Biruni_3.png} \end{center}

Re-rendered:

\begin{center} \includegraphics [scale=0.18] {BC_5.png} \end{center}

$\triangle DBZ$ is isosceles, as usual, and $DZ$ is extended to meet the circle at $H$.  To make the discussion simpler, let us refer to the measure of $\angle BDZ$ as $\alpha$ and that of $\angle DBZ$ as $\beta$.

There are several other angles all equal to $\beta$, $\angle DZB$ (by construction), $\angle AZH$ (vertical angles), $\angle AHZ$ (cuts the same arc $AD$), and $\angle GHD$ (cuts arc $GD = $ arc $AD$).

We also have the third angle in $\triangle DBZ$, namely $\angle BDZ = \alpha$ which cuts arc $BGH$.  This angle is equal to $A$ (cuts the same arc).

It follows that the sum of all the angles at $A$ and $H$ is equal to 180, from which we obtain $AB \parallel GH$.  

Parallel chords in a circle cut equal arcs and chords.

\begin{center} \includegraphics [scale=0.18] {BC_5.png} \end{center}

Thus $GB = AH$.  And since $\triangle AZH$ is isosceles, $GB = AH = AZ$.

Adding equals:

$GB + BE = AZ + ZE = AE$.

$\square$

This proof is incomplete, since for a given arrangement of $G$, $A$ and $D$, depending on the placement of $B$, one can draw $H$ in three different ways --- coincident with $G$, or lying on either side of it.

\begin{center} \includegraphics [scale=0.20] {BC_5a.png} \end{center}

The left panel is our original diagram.  In the right panel, it appears that $GH \parallel AB$.  Can the proof be rescued?  In a word, yes.

\emph{Proof}.

Just like previously, as the arc left after subtracting $GD$ and $AD$ from the whole circle, arc $AG$ must subtend $\alpha = \angle BDZ$.  So $\angle ABG = \alpha$.

\begin{center} \includegraphics [scale=0.15] {BC_5b.png} \end{center}

But $\angle BDH = \alpha$.  It follows that arc $BH = $ arc $AG$ and then by addition, arc $BG = $ arc $AH$ and so the chords are equal.  $BG = AH$.

Now proceed as in the first case.

$\square$

We might also proceed by showing that the obtuse angles at $X$, $Z$ and $B$ are all equal.  

Since the acute angles at $X$ are equal, it follows that $\triangle BDX \sim \triangle BZX \sim \triangle GHX$.

\begin{center} \includegraphics [scale=0.15] {BC_5b.png} \end{center}

We still have to show that $AH = AZ$, but this follows since $\angle AHD$ cuts arc $AD$ so $\triangle AHZ$ is isosceles.

$\square$

The \textbf{third case} is where the extension of $DZ$ terminates exactly on point $G$.  $H$ coincides with $G$.

Again, arc $AG$ corresponds to inscribed angles with measure $\alpha$ such as $\angle ABG$.

\begin{center} \includegraphics [scale=0.15] {BC_5c.png} \end{center}

But $\angle BAG$ cuts the same arc as $\angle BDZ$.

So $\triangle GAB$ is isosceles.

We have $AG = AZ$ so $AZ = BG$, and the result follows easily.

$\square$

\subsection*{sixth proof:  Suter(f)}

\emph{Proof}.

Draw the circle on center $D$ with radius $AD$.  Extend $AB$ to $Z$.  Draw $AD$ and $GD$.

\begin{center} \includegraphics [scale=0.35] {BC_6.png} \end{center}

$AZ$ is a chord of the circle on $D$.  The perpendicular $DE$ goes through the center $D$, therefore it bisects $AZ$.

We have $AE = ZE = ZB + BE$.

Three angles intercept an arc between points $A$ and $G$, $\angle BZG$ on the big circle, $\angle ABG$ on the small circle, and $\angle ADG$ on both.

Since $\angle ADG$ is a central angle in the big circle on $D$, it is twice $\angle AZG$.  But $\angle ADG = \angle ABG$.

Hence $\angle ABG$ is external to $\triangle BZG$ and twice the measure of $\angle BZG$.  It follows that $\angle BZG = \angle BGZ$.

Thus $\triangle AZG$ is isosceles with $GB = ZB$.

Substituting into what we had above:  $AE = GB + BE$.

$\square$

\subsection*{seventh proof:  Suter(p)}

\emph{Proof}.

Draw $GD$ and $AD$.  Extend $GB$ to $Z$ and connect to $D$ such that $Z$ is a right angle.

\begin{center} \includegraphics [scale=0.35] {BC_7.png} \end{center}

$\angle ZGD = \angle BAD$ because they cut the same arc.  We get the third angle by sum of angles and also $GD = AD$ so $\triangle GZD \cong \triangle AED$ by ASA.

$DZ = DE$ and then $\triangle DBZ \cong \triangle DBE$ by HL.  So $BZ = BE$.

From the first congruence $GZ = AE$.

$GZ = GB + BZ = GB + BE$.

$\square$

\subsection*{eighth proof:  Suter(n)}

As before, $AD = GD$ and $DE \perp AB$.

\begin{center} \includegraphics [scale=0.20] {BC_8.png} \end{center}

\emph{Proof}.

Extend $DE$ to cut the circle at $Z$.

Draw $AG$ and its perpendicular bisector $HTK$.  It is a diameter of the circle and goes through $D$.  Why?

Connect $ZK$ and then draw $MK \parallel DEZ$.

Since $DHTK$ is a diameter, $Z$ is right.  We're given that $\angle DEA$ is right.

Since $MK \parallel DEZ$, $EMKZ$ is a rectangle.

By similar right triangles, $D = \angle HKM = \angle A$.

So $BG = ZK = EM$.

From rectangles in a circle we know that $BE = AM$.

Add equals to equals to obtain the result.

$\square$

\subsection*{Tuan proof}

There are a couple additional proofs of the broken chord theorem that I found more recently, beyond what is in Suter.

It is on the web attributed to someone named Bùi Quang Tuån.  There is another proof from the same source of the \hyperref[sec:Pthm_Tuan]{\textbf{Pythagorean theorem}} that I like as well. 

\url{https://www.cut-the-knot.org/pythagoras/BrokenChordPythagoras.shtml}

Google also turns up a blog, but no biographical info.

\url{https://artofproblemsolving.com/community/c1598}  

This proof is based on a rectangle.

\begin{center} \includegraphics [scale=0.35] {BC_Tuan.png} \end{center}

\emph{Proof}.

Given $\angle MDA$ is right.

Find $F$ such that $MF = BG$.

Extend the diameter through $M$ to $P$.

Draw $PF$ and find $E$ on the intersection with $AB$.

Find points $Y,Z$ on the intersections with $AG$.

1. $F$ is right (Thales).

2. Since arc $AM = $ arc $GM$, $PM \perp AG$ and bisects it.

3. So $Y$ is right.

4. Also since arc $MF = $ arc $BG$, $P = A$.

5. $\triangle PYZ \sim \triangle AZE$ by (4) and vertical angles.

6. Hence $E$ is right.

\begin{center} \includegraphics [scale=0.35] {BC_Tuan.png} \end{center}

Hence $MDEF$ is a rectangle.  

We rely on the properties of parallel chords in a circle and the extension to rectangles with two vertices on the circle, as above.

Thus $BG = MF = DE$ and the main result follows easily by addition.

$\square$

from Bogomolny's site, a proof by 

\subsection*{Mariano Perez de la Cruz}

\begin{center} \includegraphics [scale=0.3] {BC_Cruz.png} \end{center}

\emph{Proof}.

Extend $AB$ to $F$ such that $AE = FE$.

Extend $DE$ to meet the circle at $Z$.  Draw $FZ$ and $AZ$.

$DEZ$ is the perpendicular bisector of $AF$.  

So $\triangle AFZ$ is isosceles and $A = F$.

By the secant-secant theorem, $\angle FGB$ is equal to $A$ and so equal to $F$.

Thus $\triangle FGB$ is isosceles, with $BG = FB$.

But $FE = BE + FB = BE + BG = AE$.

$\square$

\begin{center} \includegraphics [scale=0.3] {BC_Cruz.png} \end{center}

Note:  some additional material on the theorem is here:

\url{https://www.cut-the-knot.org/triangle/BrokenChordmpdlc.shtml}

\end{document}